% source: An algebraic approach to the Riemann-Roch theorem and the arithmetic theory of function fields (Athanasios Papaioannou), http://www.fen.bilkent.edu.tr/~franz/mat/Riemann-Roch.pdf
% pdf_page: 0006
% retrieved: 2026-07-14
% transcribed_by: page-transcriber (AI vision, no OCR)

\documentclass[11pt]{article}
\usepackage{amsmath,amssymb,amsthm}

% Small-caps theorem style matching the source (THEOREM/COROLLARY/DEFINITION headings, italic bodies).
\newtheoremstyle{scplain}{\topsep}{\topsep}{\itshape}{0pt}{\scshape}{.}{ }{\thmname{#1} \thmnote{#3}}
\theoremstyle{scplain}
\newtheorem*{theorem}{Theorem}
\newtheorem*{corollary}{Corollary}
\newtheorem*{definition}{Definition}

% Small-caps "Proof." to match the source (standard amsthm definition with \scshape head).
\makeatletter
\renewenvironment{proof}[1][\proofname]{%
  \par\pushQED{\qed}\normalfont\topsep6\p@\@plus6\p@\relax
  \trivlist\item[\hskip\labelsep\scshape #1\@addpunct{.}]\ignorespaces}%
  {\popQED\endtrivlist\@endpefalse}
\makeatother
\renewcommand{\qedsymbol}{$\square$}

% Non-standard operators used on this page (typeset upright in the source):
\DeclareMathOperator{\Div}{Div}
\DeclareMathOperator{\Supp}{Supp}
% Standard operators used: \deg, \min. Standard amssymb symbols: \mathbb, \mathbf, \varnothing.

\begin{document}

\section*{1.2\quad The Weak Approximation Theorem.}

\noindent
In this section we introduce the Weak Approximation Theorem, which will later be
used to prove that the notion of the \emph{genus} of a function field is
well-defined. The proof itself is rather technical and marginally enlightening,
so it will be omitted---it is, however, readily available in [ST], Theorem I.3.1
and Iwasawa [IW], Theorem 1.1. The precise formulation of this approximation
result has as follows

\begin{theorem}[1.9 (Weak Approximation Theorem)]
Let $\mathbb{F}/\mathbb{K}$ be a function field,
$P_1, P_2, \ldots, P_n \in \mathbf{P}_{\mathbb{F}}$ pairwise distinct places of
$\mathbb{F}/\mathbb{K}$, $x_1, x_2, \ldots, x_n \in \mathbb{F}$ and
$r_1, r_2, \ldots, r_n \in \mathbb{Z}$. Then, there is an element $x \in \mathbb{F}$
such that
\[
  v_{P_i}(x - x_i) = r_i,
\]
for $1 \leq i \leq n$.
\end{theorem}

\begin{proof}
See the above references.
\end{proof}

A different perspective explains the name given to this proposition. If we
consider the absolute value induced on $\mathbb{F}$ by the valuation $v_P$ (we
define $|z|_P = 2^{-v_P(z)}$ if $z \neq 0$ and $0$ otherwise), then the Weak
Approximation Theorem implies that the elements $x_i$ can be approximated
arbitrarily well by a single element $x \in \mathbb{F}$. From a different point of
view, the Approximation Theorem can be regarded as the topological equivalent of
the Chinese Remainder Theorem; like an approximation theorem from Number Theory,
it yields information about the solution of systems of equations.

Some direct corollaries of this fact are

\begin{corollary}[1.10]
Any function field has infinitely many places.
\end{corollary}

\begin{proof}
Assume that there only finitely many places $P_1, P_2, \ldots, P_n$; by the above
theorem, we can find a non-zero element $x \in \mathbb{F}$ such that
$v_{P_i}(x) > 0$, for $1 \leq i \leq n$. Then $x$ is transcendental over
$\mathbb{K}$ but has no poles, a contradiction to .
\end{proof}

\begin{theorem}[1.11]
Let $\mathbb{F}/\mathbb{K}$ be a function field and $P_1, P_2, \ldots, P_r$be
zeroes of a fixed element $x \in \mathbb{F}$. Then
\[
  \sum_{i=1}^{r} v_{P_i}(x) \deg P_i \leq [\mathbb{F} : \mathbb{K}(x_i)].
\]
\end{theorem}

\begin{proof}
See the references made in the introduction.
\end{proof}

\begin{corollary}[1.12]
Any element $x \in \mathbb{F}^*$ has finitely many places and zeroes.
\end{corollary}

\begin{proof}
If $x$ is constant (namely, non transcendental), then it has neither zeroes nor
poles, and the claim holds trivially. Otherwise, the number of zeroes of $x$ is
bounded above by $[\mathbb{F} : \mathbb{K}(x)]$ by the previous theorem; the same
argument shows that the number of zeroes of $x^{-1}$ is finite.
\end{proof}

\section*{1.3\quad Divisors.}

\noindent
Since the closure $\tilde{K}$ of the base field in an algebraic function field
$\mathbb{F}/\mathbb{K}$ is a finite extension over $\mathbb{K}$,
$\mathbb{F}/\mathbb{K}$ is also a function field, and in this case the base field
is the full constant field of the function field; this situation offers certain
technical advantages over the more general one, and henceforth we will assume
that the base field is algebraically closed in the extension $\mathbb{F}/\mathbb{K}$.

The basic notion of this section is that of a divisor; it is given by the
following

\begin{definition}[1.6]
Let $\Div(\mathbb{F})$ denote the free abelian group generated by the places of
$\mathbb{F}/\mathbb{K}$. The elements $D$ of $\Div(\mathbb{F})$ are called divisors
and can be presented as formal sums
\[
  D = \sum_{P \in \mathbf{P}_{\mathbb{F}}} n_P P,
\]
where the coefficients $n_p \in \mathbb{Z}$ are $0$ for almost all
$P \in \mathbf{P}_{\mathbb{F}}$. We also define the support $\Supp(D)$ of $D$ to be
the set of all $P \in \mathbf{P}_{\mathbb{F}}$ such that $n_P \neq 0$ in the above
representation of $D$.
\end{definition}

\begin{center}6\end{center}

\end{document}
